\chapter{La puberté}\label{ch:g8:puberty}

Quelque part entre neuf et seize ans environ, tout corps humain se
reconstruit~: la seconde poussée de croissance de
\cref{rem:g3:stages-of-life:bursts}, et bien plus que la taille.
\Cref{rem:g5:human-reproduction:adults} annonçait ce chapitre il y a
des années. Le voici --- les faits, la fameuse injustice du
calendrier, et la machinerie, qui va nous offrir une découverte plus
grande que la \omterm{def:g8:puberty:puberty}{puberté} elle-même~: le corps possède un second système
de communication.

\section{Ce qui change}

\begin{definition}[Puberté]\label{def:g8:puberty:puberty}
La \emph{puberté}\index{puberté} est l'étape au cours de laquelle un
corps d'enfant devient un corps d'adulte --- avant tout un corps
d'adulte au sens de \cref{ch:g8:sexual-reproduction}~: les organes
reproducteurs arrivent à maturité et commencent à produire des
\emph{\omterm{def:g8:sexual-reproduction:gametes}{gamètes}}, et le corps prend sa forme adulte. C'est le \omterm{def:g4:heart-and-blood:heart}{cœur}
biologique de l'adolescence
(\cref{def:g3:stages-of-life:stages}).
\end{definition}

\begin{proposition}[Les changements visibles]\label{prop:g8:puberty:changes}
En plus de la poussée de croissance, sur plusieurs années~:
\begin{itemize}
  \item chez les filles, couramment de neuf à treize ans environ~:
        les seins se développent, les hanches s'élargissent, une
        pilosité apparaît --- et les \emph{premières règles}
        marquent la mise en route des cycles de
        \cref{ch:g8:reproductive-systems}~;
  \item chez les garçons, couramment un an ou deux plus tard~: la
        voix devient plus grave (elle «~mue~»), les \omterm{def:g1:my-body:joint}{épaules}
        s'élargissent, une pilosité corporelle et faciale apparaît,
        et les testicules commencent à produire des \omterm{def:g4:animal-reproduction:sexual}{cellules
        mâles}~;
  \item chez les deux~: la \omterm{def:g1:the-five-senses:sense}{peau} devient plus grasse (l'âge des
        boutons), la \omterm{rem:g7:removing-wastes:sweat}{sueur} change, l'humeur suit la reconstruction
        --- le \omterm{def:g5:brain-in-command:system}{cerveau} aussi est recâblé vers sa forme adulte.
\end{itemize}
\end{proposition}
\admitted

\begin{remark}[Le calendrier est personnel]\label{rem:g8:puberty:timetable}
\Cref{rem:g1:my-body:different} le disait pour la taille~; la
\omterm{def:g8:puberty:puberty}{puberté} le répète plus fort. Des camarades du même âge peuvent se
trouver à des années d'écart sur cette route --- précoce et tardif
sont normaux tous les deux, l'ordre des changements varie, et
presque tout le monde arrive au même endroit au bout du compte. Les
intervalles ci-dessus sont larges, et le plus large en eux est le
peu qu'ils disent d'une personne en particulier.
\end{remark}

\section{La machinerie~: un second système de messagers}

\begin{definition}[Les hormones, première rencontre]\label{def:g8:puberty:hormones}
La \omterm{def:g8:puberty:puberty}{puberté} ne se décide pas dans les organes qu'elle change. Elle
est \emph{commandée} --- par des messagers chimiques appelés
\emph{hormones}\index{hormone}~: des substances libérées dans le
\emph{\omterm{prop:g4:journey-of-food:absorb}{sang}} par des organes spécialisés, transportées partout, et
obéies par les organes bâtis pour les entendre. À la \omterm{def:g8:puberty:puberty}{puberté}, un
centre de commande à la base du \omterm{def:g5:brain-in-command:system}{cerveau} se réveille, ses hormones
donnent leurs ordres aux organes reproducteurs, et \emph{leurs}
hormones --- portées dans tout le corps par le service de livraison
de \cref{prop:g4:heart-and-blood:carries} --- commandent toute la
reconstruction~: croissance, pilosité, voix, cycles.
\end{definition}

\begin{example}[Un seul signal, beaucoup de chantiers]\label{ex:g8:puberty:sites}
Le \omterm{prop:g4:journey-of-food:absorb}{sang} porte une \omterm{def:g8:puberty:hormones}{hormone} partout à la fois --- et c'est pourquoi la
\omterm{def:g8:puberty:puberty}{puberté} est une reconstruction \emph{coordonnée}~: la voix, les
\omterm{def:g1:my-body:joint}{épaules} et la pilosité suivent le même messager, sur un seul
calendrier, alors que les chantiers sont à une demi-longueur de
corps les uns des autres. Aucun câblage nerveux ne relie un larynx à
une barbe~; la diffusion par le \omterm{prop:g4:journey-of-food:absorb}{sang}, si
(\cref{ch:g8:hormonal-communication} fait de ce système un chapitre
à part entière).
\end{example}

\begin{remark}[Pourquoi cette météo d'humeur~?]\label{rem:g8:puberty:moods}
Les mêmes messagers atteignent le \omterm{def:g5:brain-in-command:system}{cerveau}, lui-même en pleine
rénovation~: sautes d'humeur, intensités nouvelles, \omterm{prop:g1:healthy-habits:needs}{sommeil} qui
glisse vers plus tard --- de la météo, pas une faute. Connaître la
machinerie aide des deux côtés~: les orages sont de la chimie plus
un chantier, ils sont normaux, et ils passent.
\end{remark}

\section{Vivre pendant la reconstruction}

\begin{proposition}[Prendre soin de soi pendant le chantier]\label{prop:g8:puberty:care}
Un corps en pleine reconstruction a des besoins de bâtisseur, tous
déjà rencontrés dans ce livre~:
\begin{itemize}
  \item du matériau et du carburant~: la poussée se paie en matière
        --- une alimentation équilibrée
        (\cref{met:g3:balanced-diet:check}), matière de
        construction comprise
        (\cref{exo:g7:digestion-and-nutrients:11})~;
  \item du \omterm{prop:g1:healthy-habits:needs}{sommeil}~: la rénovation du \omterm{def:g5:brain-in-command:system}{cerveau} et la croissance du
        corps se font toutes deux la nuit
        (\cref{prop:g5:brain-in-command:protect}) --- et l'horloge
        qui glisse vers plus tard a toujours besoin de ses heures~;
  \item du mouvement~: les \omterm{def:g3:bones-and-muscles:skeleton}{os} et le \omterm{def:g4:heart-and-blood:heart}{cœur} sont dans leurs dernières
        et meilleures années de construction
        (\cref{met:g3:bones-and-muscles:care})~;
  \item et aucun sabotage~: les substances de
        \cref{prop:g5:healthy-choices:substances} frappent le plus
        durement les corps et les cerveaux encore en chantier.
\end{itemize}
\end{proposition}
\admitted

\begin{example}[Les questions ont des adresses]\label{ex:g8:puberty:questions}
La \omterm{def:g8:puberty:puberty}{puberté} soulève des questions --- sur le calendrier, sur les
corps, sur les sentiments --- et elles ont des adresses en règle~:
les infirmières scolaires et les médecins y répondent chaque jour,
en confidence et sans s'étonner~; les adultes de confiance sont tous
passés par là~; et les chapitres de ce livre sont faits pour être
relus. Ce qui ne mérite aucune confiance~: la légende de la cour de
récréation et les pires recoins d'internet, qui font commerce des
incertitudes exactes de cet âge.
\end{example}

\section{Exercices}

\begin{exercise}[$\star$]\label{exo:g8:puberty:1}
Définis la \omterm{def:g8:puberty:puberty}{puberté}, et nomme l'étape de
\cref{def:g3:stages-of-life:stages} dont elle est le \omterm{def:g4:heart-and-blood:heart}{cœur}.
\end{exercise}

\begin{exercise}[$\star$]\label{exo:g8:puberty:2}
Énumère trois changements de la \omterm{def:g8:puberty:puberty}{puberté} des filles, trois de celle
des garçons, et deux communs aux deux.
\end{exercise}

\begin{exercise}[$\star$]\label{exo:g8:puberty:3}
Qu'est-ce qui marque le début de la production de \omterm{def:g8:sexual-reproduction:gametes}{gamètes} chez
chaque sexe, d'après \cref{prop:g8:puberty:changes}~?
\end{exercise}

\begin{exercise}[$\star$]\label{exo:g8:puberty:4}
Qu'est-ce qu'une \omterm{def:g8:puberty:hormones}{hormone} --- libérée où, transportée comment, obéie
par qui~?
\end{exercise}

\begin{exercise}[$\star$]\label{exo:g8:puberty:5}
Où commence la chaîne de commande de la \omterm{def:g8:puberty:puberty}{puberté}, et où finit-elle~?
\end{exercise}

\begin{exercise}[$\star$]\label{exo:g8:puberty:6}
Deux jeunes de treize ans sont visiblement à des stades différents.
Que dit \cref{rem:g8:puberty:timetable} --- entièrement~?
\end{exercise}

\begin{exercise}[$\star\star$]\label{exo:g8:puberty:7}
Pourquoi un seul messager peut-il coordonner des chantiers séparés
par une demi-longueur de corps, là où des \omterm{def:g5:brain-in-command:system}{nerfs} demanderaient un
câblage jusqu'à chacun~? Sers-toi de \cref{ex:g8:puberty:sites}.
\end{exercise}

\begin{exercise}[$\star\star$]\label{exo:g8:puberty:8}
Explique la météo d'humeur de \cref{rem:g8:puberty:moods} --- deux
causes, une réassurance.
\end{exercise}

\begin{exercise}[$\star\star$]\label{exo:g8:puberty:9}
Relie la poussée de croissance à l'assiette~: quelles familles de
\omterm{def:g7:digestion-and-nutrients:families}{nutriments} la reconstruction sollicite-t-elle le plus, et d'après le
raisonnement de quel chapitre~?
\end{exercise}

\begin{exercise}[$\star\star$]\label{exo:g8:puberty:10}
Pourquoi les substances de
\cref{prop:g5:healthy-choices:substances} frappent-elles le plus
durement un corps d'\omterm{def:g3:stages-of-life:stages}{adolescent}~? (L'image du chantier porte la
réponse.)
\end{exercise}

\begin{exercise}[$\star\star$]\label{exo:g8:puberty:11}
Classe en fiable et non fiable, avec une raison pour chacun~:
l'infirmière scolaire~; les promesses miracles d'une vidéo virale~;
ce livre~; la légende de la cour de récréation~; le médecin de
famille.
\end{exercise}

\begin{exercise}[$\star\star\star$]\label{exo:g8:puberty:12}
«~La \omterm{def:g8:puberty:puberty}{puberté} est l'entrée en scène publique du système
hormonal.~» Défends la phrase~: que la reconstruction
démontre-t-elle sur la portée de ce système, sa coordination et sa
différence avec la signalisation nerveuse --- trois affirmations,
trois observations.
\end{exercise}

\section{Problème~: la réunion sur les courbes de croissance}

\begin{problem}[{Problème du week-end --- quatre courbes de
croissance (imaginaires), lues avec tout le chapitre}]\label{pb:g8:puberty:1}
Un médecin scolaire examine quatre courbes, taille en fonction de
l'âge, chacune avec ses dates de mesure~: L, dont la courbe a bondi
de onze à treize ans et s'aplatit maintenant à quinze~; M, en montée
régulière à treize ans, poussée pas encore commencée~; N, en pleine
poussée à quatorze~; O, dont le segment raide a couru de neuf à
douze ans et qui a atteint un plateau.

\medskip\noindent\textbf{Partie I --- Lire les courbes.}
\begin{enumerate}
  \item Place chaque courbe sur la carte de
        \cref{rem:g3:stages-of-life:bursts}~: quelle poussée
        apparaît, et où chaque élève se situe-t-il dessus~?
  \item M s'inquiète d'être «~à la traîne~». Rédige la réponse du
        médecin à partir de
        \cref{rem:g8:puberty:timetable}, avec la prédiction que
        soutient la montée régulière de sa courbe.
  \item O demande si le plateau signifie «~fini de grandir à douze
        ans~». Corrige cette lecture~: qu'est-ce que le plateau
        termine, et qu'est-ce qui continue après lui~?
  \item Pourquoi le médecin lit-il la \emph{pente}, et non la
        taille, pour juger où chacun se situe~?
\end{enumerate}

\medskip\noindent\textbf{Partie II --- Derrière les courbes.}
\begin{enumerate}[resume]
  \item Nomme la machinerie qui a déclenché chaque poussée, et le
        trajet qu'ont suivi ses ordres.
  \item Le médecin suit aussi le \omterm{prop:g1:healthy-habits:needs}{sommeil} et l'alimentation à ces
        consultations. Relie chacun à sa ligne de
        \cref{prop:g8:puberty:care}.
  \item N signale des orages d'humeur et un coucher qui glisse vers
        plus tard. Donne l'explication en deux parties de
        \cref{rem:g8:puberty:moods} --- et la conséquence pratique
        pour le \omterm{prop:g1:healthy-habits:needs}{sommeil} de N les jours d'école.
  \item L pose au médecin, en privé, une question sur un changement
        que la courbe ne peut pas montrer. Dis, sans inventer la
        question, pourquoi le cabinet est la bonne adresse
        (\cref{ex:g8:puberty:questions}).
\end{enumerate}

\medskip\noindent\textbf{Partie III --- La note de
réunion.}
\begin{enumerate}[resume]
  \item Rédige la note d'une ligne pour chaque courbe~: stade,
        normalité, et point de vigilance éventuel.
  \item Les quatre courbes diffèrent de plusieurs années, et
        pourtant le médecin les classe toutes en «~normal~».
        Justifie en une phrase.
  \item Un parent demande «~quelque chose pour accélérer M~».
        Réponds avec la remarque sur le calendrier et avec la
        machinerie~: qu'y a-t-il à corriger~?
  \item Termine par le résumé du chapitre en une phrase pour le
        panneau d'affichage de la classe~: ce qu'est la \omterm{def:g8:puberty:puberty}{puberté},
        qui la dirige, et ce dont elle a besoin.
\end{enumerate}
\end{problem}
